\section{Problem Setting and Threat Model}
\subsection{Communication Setting}
We consider a sender-receiver pair and a passive observer. Alice (sender) and Bob (receiver) share a prompt and a password. Alice encrypts a secret message into a packet and embeds that packet into generated text. Bob receives only the text and attempts recovery using the same prompt and password. Censor observes the transmitted text but does not participate in key exchange or protocol synchronization.

\subsection{Adversary Scope}
Our baseline adversary is a passive Censor who can read the transmitted text but does not know the shared prompt or password and does not actively modify the text in transit. We do not assume token insertions, deletions, rewrites, or reordering by the channel in this baseline model. Under this scope, the target security statement is conservative: stego text should be indistinguishable from natural text to the passive observer to the extent supported by our assumptions and evidence.

\subsection{Operational Assumptions}
Deterministic recovery in Ghostext requires strict runtime synchronization between Alice and Bob. In particular, model/backend identity, tokenizer behavior, prompt content, prompt-template behavior, seed, candidate policy parameters, quantization settings, and interval-code parameters must match. These assumptions are explicit and are part of the protocol surface.

\subsection{Security Objectives}
Within this model, Ghostext targets three objectives. First, recoverability: when configuration is matched, Bob should recover the exact plaintext. Second, fail-closed behavior: when synchronization breaks, decoding should return an explicit error instead of an approximate plaintext. Third, reproducibility: protocol behavior and reported metrics should be rerunnable from recorded configuration and scripts.

\subsection{Non-Goals in This Version}
This version does not claim robustness under active text edits, resistance against dedicated steganalysis detectors, or cross-model and cross-version decode compatibility. These are important directions but currently outside the supported guarantee boundary.
